

%% ===================================================================
\section{\vrev{이중채널 감쇠 메커니즘}}
\label{sec:ch5_dual_channel}
%% ===================================================================

C2--C3 비교에서 Cohen's $d = -0.238$라는 평균 수준의 비차별성과,
C3의 표준편차가 C2의 $1.29$배라는 분산 수준의 차이 사이에는
\jrev{구조적 패러독스(structural paradox)}가 존재한다. 이 절에서는 이 \jrev{패러독스}의 원인을 P85.8 교차점에서의
분포 역전과 이중채널 감쇠 메커니즘으로 해명한다.


%% -------------------------------------------------------------------
\subsection{\vrev{P85.8 교차점(분포 역전의 전환점)}}
\label{subsec:ch5_crossover}
%% -------------------------------------------------------------------

C2와 C3의 $\eta$ \jrev{백분위 함수가 교차하는 지점\chg{(}$\text{VaR}_p^{\text{C2}} = \text{VaR}_p^{\text{C3}}$를}
만족하는 백분위수 $p$\chg{)}은 약 \textbf{P85.8}에서 관측된다.
이 교차점 이하에서는 \jrev{$\text{VaR}_p^{\text{C2}} > \text{VaR}_p^{\text{C3}}$}, 즉 C2가 C3보다
높은 $\eta$를 보이며, 이는 가우시안 분포의 중심 집중이 파레토보다
높은 중앙값을 산출하기 때문이다.
교차점 이상에서는 \jrev{$\text{VaR}_p^{\text{C3}} > \text{VaR}_p^{\text{C2}}$}로 역전되어,
파레토 효과가 급격히 발현한다.

이 구조적 역전의 실무적 함의는 중대하다.
전체 표본의 $85.8\%$에 해당하는 일상적 운영 상황에서는 C2와 C3의 차이가 미미하여
조직은 두 불확실성 유형을 구분할 실무적 필요가 없다.
그러나 상위 $14.2\%$의 극단 상황에서는 파레토 불확실성이 복구 시간을
급격히 연장한다. $\mathrm{CVaR}_{99}$ 기준 C3가 C2 대비 $+24.0\%$ 높은 것은
이 극단 영역의 격차를 정량화한 것이다.

이러한 ``평균에서의 무차별$+$극단에서의 극적 분기'' 패턴은
BCM 실무에서 극단 사건에 대한 별도의 비상 계획\chg{(}CVaR 기반의
꼬리 리스크 관리\chg{)}이 필요하다는 이론적 근거를 제공한다.


%% -------------------------------------------------------------------
\subsection{\vrev{분할 회귀와 $k_O$ 계수 감쇠}}
\label{subsec:ch5_split_regression}
%% -------------------------------------------------------------------

P85.8 교차점을 기준으로 C3의 표본을 Core($\leq$ P85.8, $n = 8{,}580$)와
Tail($>$ P85.8, $n = 1{,}420$)로 분할하여 회귀 분석을 수행하면,
관측성 가중치 $k_O$의 효과가 두 영역에서 질적으로 상이한 구조를 보인다.

Core 영역의 1차 단순 회귀에서 $k_O$ 계수는 $-0.798$이다.
$k_O$가 1단위 증가하면 $\eta$가 $0.798$만큼 감소하며,
이는 $\DRTO$가 $0.798 \times R_0 = 4.79$시간 단축됨을 의미한다.
Tail 영역의 1차 단순 회귀에서 $k_O$ 계수는 $-0.415$이다.
동일한 $k_O$ 1단위 증가가 Tail에서는 $\eta$를 $0.415$만큼 감소시켜,
$\DRTO$가 $0.415 \times R_0 = 2.49$시간 단축된다.

Tail 대비 Core의 계수 비율은 $|-0.415|/|-0.798| \approx 0.52$배이다.
이 $0.52$배 감쇠가 \textbf{이중채널 감쇠(dual-channel attenuation)}의
핵심 지표이며, 관측성 가중치와 파레토 불확실성이 결합될 때
극단 영역에서 관측성 효과가 상대적으로 약화되는 구조적 증거이다.


%% -------------------------------------------------------------------
\subsection{\vrev{이중채널 감쇠의 구조적 원인}}
\label{subsec:ch5_mechanism}
%% -------------------------------------------------------------------

이중채널 감쇠의 구조적 원인은 개별 시그모이드 바운딩 함수의 비선형성에 있다.

\chg{먼저, Core 채널은 선형 영역에서 작동한다.}
Core 영역($\leq$ P85.8)에서는 불확실성 $U_{\text{raw}}$의 값이 상대적으로 작아
$z_U = \kappa \cdot k_U \cdot R_0 \cdot U_{\text{raw}} / (R_{\max} - R_0)$이 시그모이드의 선형 근사 영역($|z| \ll 1$)에
위치한다. 이 영역에서 $S(z) \approx z$이므로,
$E_{U,\text{bnd}} \approx (R_{\max} - R_0) \cdot z_U$로 불확실성의 효과가
입력에 거의 비례한다. $k_O$의 효과는 $E_{O,\text{bnd}}$를 통해
단조적으로 전달되며, 관측성의 복구 단축 효과가 충분히 발현된다.

\chg{반면에, Tail 채널은 비선형 구간에서 포화 효과를 겪는다.}
Tail 영역($>$ P85.8)에서는 큰 $U_{\text{raw}}$ 값(파레토 분포의 우측 꼬리)에 의해
$z_U$가 시그모이드의 곡률(curvature)이 큰 영역에 진입한다.
이 상태에서 $E_{U,\text{bnd}}$가 시그모이드의 포화에 근접하여
불확실성 채널이 $\DRTO$를 지배하게 되며,
$k_O$의 변화가 $E_{O,\text{bnd}}$를 통해 전달되는 관측성의
복구 단축 효과가 불확실성의 상향 효과에 의해 상대적으로 약화된다.
구체적으로, $\DRTO = R_0 + E_{O,\text{bnd}} + E_{U,\text{bnd}}$에서
$E_{U,\text{bnd}}$가 지배적일수록 $E_{O,\text{bnd}}$의 상대적 기여가 감소한다.

정리하면, 파레토 불확실성은 시그모이드의 비선형 구간을 활성화하여
불확실성 채널을 지배적으로 만들며, 이로 인해 관측성의 상대적 역할이 축소된다.
이는 극단 불확실성 환경에서 관측성 투자만으로는 복구 시간 단축의
효과가 제한적임을 의미하며, 불확실성 자체의 관리가 우선되어야 한다는
실무적 함의로 직결된다.


%% -------------------------------------------------------------------
\subsection{\vrev{분포별 회귀 구조 차이(지배 변수의 전환)}}
\label{subsec:ch5_dominant_variable}
%% -------------------------------------------------------------------

4.4(회귀분석 결과)의 분석에서 C2와 C3의 변수 진입 순서가 질적으로
다르다는 발견은 교차 분석 관점에서 중요한 의미를 갖는다.
\jrev{\tabref{tab:ch5_dominant_var}\refpo{tab:ch5_dominant_var}{은}{는}} 두 조건의 변수 진입 1순위 및
2순위까지의 누적 $R^2$를 비교한 것이다.

\begin{table}[htbp]
  \centering
  \caption{\jrev{C2 vs.\ C3 전진 선택법 변수 진입 구조 비교}}
  \label{tab:ch5_dominant_var}
  \begin{tabular}{l cc cc}
    \toprule
    & \multicolumn{2}{c}{\textbf{C2 (Gaussian)}} & \multicolumn{2}{c}{\textbf{C3 (Pareto)}} \\
    \cmidrule(lr){2-3} \cmidrule(lr){4-5}
    진입 순서 & 변수 & 누적 $R^2$ & 변수 & 누적 $R^2$ \\
    \midrule
    1순위 & $k_O$ & 0.768 & $U$ & 0.424 \\
    2순위 & $U$ & 0.931 & $k_O$ & 0.686 \\
    \bottomrule
  \end{tabular}
\end{table}

\chg{지배 변수 전환의 의미를 보면,}
C2에서 $k_O$가 1순위($R^2 = 0.768$)인 것은 가우시안 불확실성의
분산이 상대적으로 작아($\mathrm{SD} = 1.0$) $\eta$ 변동의 주된 원천이
관측성임을 반영한다. 반면, C3에서 $U$가 1순위($R^2 = 0.424$)인 것은
파레토 분포의 높은 분산(\jrev{$U$ 값의 이론적 분산이 가우시안의 $27$배})이
$\eta$ 변동에서 불확실성의 역할을 증대시키기 때문이다.

\chg{2순위 투입 후 설명력 격차를 보면,}
C2에서 2순위($U$) 투입 후 $R^2 = 0.931$에 도달하는 반면,
C3에서 2순위($k_O$) 투입 후 $R^2 = 0.686$에 머문다.
이 $0.245$의 격차는 C3의 $\eta$ 변동 중 상당 부분이
비선형 교호 작용에 의해 설명됨을 시사하며,
\pdferr{2차 회귀} 모형의 필요성을 교차 분석 관점에서 뒷받침한다.

\chg{실무적 함의로서,}
이러한 지배 변수 전환은 BCM 실무에서 중요한 의사결정 원칙을 시사한다.
가우시안 불확실성(일상적 변동) 환경에서는 \textbf{관측성 투자가
최우선 전략}인 반면, 파레토 불확실성(극단 사건 가능성) 환경에서는
\textbf{불확실성 정량화가 최우선 전략}이다.
두 전략의 우선순위는 조직이 직면하는 불확실성의 분포 형태에 의해 결정되며,
이는 산업 특성에 따라 차별화되어야 한다.


\jrev{\figref{fig:ch5_dual_channel_mechanism}\refpo{fig:ch5_dual_channel_mechanism}{은}{는}} 이중 채널 감쇠의
메커니즘을 도식으로 나타낸 것이다.
\chg{그림은 P85.8 기준으로 분할한 후 Tail 채널에서 $k_O$ 회귀 계수가 Core 대비 $0.52$배로 감쇠함을 보인다.}

\begin{figure}[htbp]
  \centering
  \resizebox{\textwidth}{!}{%
  \begin{tikzpicture}[
    channel/.style={draw, rounded corners=5pt, minimum width=52mm, minimum height=14mm,
      align=center, font=\small, line width=0.6pt},
    input/.style={draw, rounded corners=4pt, minimum width=38mm, minimum height=11mm,
      align=center, font=\small, line width=0.6pt},
    result/.style={draw, rounded corners=5pt, minimum width=55mm, minimum height=14mm,
      align=center, font=\small\bfseries, line width=0.6pt},
    arr/.style={-{Stealth[length=2.5mm]}, thick},
    splitlbl/.style={font=\scriptsize, fill=white, inner sep=2pt,
      rounded corners=1pt},
    ]

    % Input variables
    \node[input, fill=blue!15] (ko) at (0, 4)   {관측성 가중치\\$k_O \in [0,1]$};
    \node[input, fill=green!15] (O) at (0, 2.2)  {관측성 원시 점수\\$O_{\text{raw}} \in [0,1]$};
    \node[input, fill=orange!15] (U) at (0, 0.4) {불확실성\\$U_{\text{raw}}^{(P)} \sim 6 \cdot \mathrm{Pareto}(\alpha\!=\!3)$};

    % Sigmoid bounding
    \node[channel, fill=gray!15] (sig) at (7, 2.2) {개별 시그모이드 바운딩\\$\DRTO = R_0 + E_{O,\text{bnd}} + E_{U,\text{bnd}}$};

    % Fork point
    \coordinate (fork) at (10.0, 2.2);

    % Split channels
    \node[channel, fill=blue!10] (core) at (14.5, 4.2)
      {\textbf{Core 채널} ($\leq$ P85.8)\\$k_O$ 계수 $= -0.798$,\; $R^2 = 0.999$};
    \node[channel, fill=red!10] (tail) at (14.5, 0.2)
      {\textbf{Tail 채널} ($>$ P85.8)\\$k_O$ 계수 $= -0.415$,\; $R^2 = 0.710$};

    % Attenuation result
    \node[result, fill=yellow!15] (amp) at (14.5, -2.2)
      {이중채널 감쇠 비율\\$|-0.415| / |-0.798| \approx 0.52$배};

    % Input arrows — 3개 모두 바운딩 박스 좌측면(west)으로 수렴
    \draw[arr, blue!60, line width=1.2pt]
      (ko.east) -- ++(1.8, 0) |- (sig.170);
    \draw[arr, green!50!black, line width=1.2pt]
      (O.east) -- (sig.west);
    \draw[arr, orange!60, line width=1.2pt]
      (U.east) -- ++(1.8, 0) |- (sig.190);
    % Arrow labels
    \node[font=\tiny, fill=white, inner sep=1pt, text=black] at (3.5, 3.3) {관측성 ($k_O$)};
    \node[font=\tiny, fill=white, inner sep=1pt, text=black] at (3.5, 2.2) {관측성 ($O$)};
    \node[font=\tiny, fill=white, inner sep=1pt, text=black] at (3.5, 1.1) {불확실성 ($U$)};

    % Sigmoid to fork
    \draw[arr] (sig.east) -- (fork);
    \draw[arr] (fork) |- (core.west);
    \draw[arr] (fork) |- (tail.west);

    % P85.8 label
    \node[splitlbl] at (10.0, 3.3) {P85.8 기준};
    \node[splitlbl] at (10.0, 2.85) {분할};

    % Tail to amplification (직선)
    \draw[arr, dashed, red!60, line width=1.2pt] (tail.south) -- (amp.north);
    % Core to amplification (우측 우회: core → 오른쪽 → amp 우측)
    \draw[arr, dashed, blue!40]
      (core.east) -- ++(1.2, 0) |- (amp.east);

    % Key insight annotation
    \node[font=\scriptsize, text=black, align=center, anchor=west] at (16.5, -0.8)
      {Tail에서 $k_O$\\$0.52$배 감쇠};
  \end{tikzpicture}
  }% end resizebox
  \caption{\jrev{이중채널 감쇠 메커니즘}}
  \label{fig:ch5_dual_channel_mechanism}
\end{figure}


%% -------------------------------------------------------------------
\subsection{\vrev{이중 채널 감쇠의 실무적 함의}}
\label{subsec:ch5_dual_practical}
%% -------------------------------------------------------------------

이중채널 감쇠 효과의 실무적 함의는 다음의 세 가지로 정리된다.

\chg{첫째, 관측성 투자의 조건부 수익 구조이다. 일상 운영(Core)에서 관측성 가중치 $k_O$를 $0.5$에서
$0.6$으로 $0.1$만큼 개선하면 $\eta$가 $0.080$만큼 감소하며($0.798 \times 0.1 = 0.080$, $\DRTO$로 환산하면
$0.48$시간), 극단 상황(Tail)에서 동일한 $k_O$ 개선은 $\eta$를 $0.042$만큼 감소시킨다($0.415 \times 0.1 = 0.042$,
$\DRTO$로 환산하면 $0.25$시간). 즉 극단 상황에서 관측성 투자의 한계 수익은 일상 운영의 $0.52$배로
감쇠된다. 둘째, 이중축 BCM 전략의 근거이다. P85.8 교차점은 BCM 전략을 일상 운영 축(C2 기반)과 극단
시나리오 축(C3 기반)으로 분리하는 자연적 경계점이며, $\mathrm{CVaR}_{99}$ 기준 $+24.0\%$의 격차는 극단
시나리오에 대비한 추가 자원 확보의 정량적 근거가 된다. 셋째, 스트레스 테스트 설계 기준이다. $0.52$배
감쇠 효과는 BCM 스트레스 테스트에서 Tail 영역의 시나리오를 별도로 설계해야 함을 의미하며, 극단
상황에서 관측성 투자의 효과가 감소하므로 불확실성 관리 자체에 초점을 맞춘 극단 시나리오 대응 전략이
필수적이다.}

\chg{H8의 검증 결론을 정리하면,} \mrev{H8은 ``Tail ($>$P85.8) 영역에서 $k_O$ 계수의 절댓값이 Core 영역의
$\leq 1.0$배''를
기준으로 한다. Tail $|\beta_{k_O}|
= 0.415$ vs.\ Core $|\beta_{k_O}| = 0.798$의 비율이 $0.415/0.798 = 0.52
\leq 1.0$이므로, \textbf{H8은 채택된다}. 이 감쇠 효과는 분포 환경에 의존하는
관측성 효과의 비선형성을 정량적으로 입증한 본 연구의 독자적
발견이다.}
